% ECONOMICS_PROBLEM: {"id": "D63-448", "title": "Fair chores with unequal time requirements", "jel_code": "D63", "source_id": "OP-0277"}
% !TeX program = xelatex
\documentclass[11pt,a4paper]{article}
\usepackage[margin=23mm]{geometry}
\usepackage{fontspec}
\setmainfont{Latin Modern Roman}
\usepackage{amsmath,amssymb,mathtools}
\usepackage{xcolor,enumitem,booktabs,longtable,array,xurl,hyperref}
\definecolor{muted}{HTML}{53636C}
\hypersetup{colorlinks=true,urlcolor=blue,linkcolor=blue}
\setlength{\parindent}{0pt}
\setlength{\parskip}{5pt}
\newcommand{\R}{\mathbb R}
\newcommand{\E}{\mathbb E}
\newcommand{\N}{\mathbb N}
\newcommand{\ind}{\mathbf 1}
\newcommand{\Var}{\operatorname{Var}}
\newcommand{\Cov}{\operatorname{Cov}}
\newcommand{\supp}{\operatorname{supp}}
\DeclareMathOperator*{\argmin}{arg\,min}
\DeclareMathOperator*{\argmax}{arg\,max}
\newcommand{\entrymeta}[3]{{\sffamily\small\color{muted}\textbf{\texttt{#1}}\quad|\quad #2\par #3\par}}
\newcommand{\Problemref}[1]{\texttt{#1}}
\newcommand{\JELref}[2]{\texttt{#2} (JEL \texttt{#1})}
\begin{document}
\section*{Fair chores with unequal time requirements}
\textbf{Problem ID:} \texttt{D63-448}\quad \textbf{JEL:} \texttt{D63}\par
% BEGIN ECONOMICS STATEMENT
\label{problem:OP-0277}
{\sffamily\small\color{muted}Original subject group: Economic theory, equilibrium and social choice\par}
\label{definitions:problem:30007559}
\textbf{Audit classification:} Published research agenda. Dagstuhl §4.1 explicitly asks for suitable envy definitions with agent-dependent chore times. The displayed approximation requirement is editorial, not the paper’s adopted notion.
\subsection*{Definitions and precise research target}
\textbf{Model and research target.} Let \(N=\{1,\ldots,n\}\) be agents and \(C\) a finite set of indivisible chores. Agent \(i\) has additive nonnegative disutility \(d_i\), person-dependent processing times \(t_{ic}>0\), and time budget \(B_i>0\). A complete assignment \(A=(A_i)_{i\in N}\) partitions \(C\) and is feasible when \(\sum_{c\in A_i}t_{ic}\leq B_i\) for every \(i\). Restrict inputs to those admitting some feasible complete assignment. For comparison, define a proposed budget-aware envy factor \(\alpha\geq1\): whenever \(i\ne j\) and \(A_j\) would fit agent \(i\)'s budget, require either \(A_i=\varnothing\), or some \(c\in A_i\) satisfying
\[
 d_i(A_i\setminus\{c\})\leq\alpha d_i(A_j).
\]
Determine the smallest factor guaranteed over a specified domain of processing times, budgets, and disutilities, and characterize when the factor-one requirement is compatible with feasible completeness. Produce tight counterexamples when budgets prevent any finite guarantee, and identify minimally restrictive conditions restoring useful bounds. Evaluate whether the comparison above has an adequate welfare interpretation when agent-specific processing times differ, or propose and axiomatize a more appropriate comparison.

\emph{Definitions and maintained assumptions (editorial unless a source definition is named).} For each \(S\subseteq C\), set \(d_i(S)=\sum_{c\in S}d_i(c)\) with \(d_i(c)\ge0\), and \(t_i(S)=\sum_{c\in S}t_{ic}\). Feasibility is \(t_i(A_i)\le B_i\) for every assigned agent. The proposed budget-aware EF1 condition compares \(i\) with \(j\) only when \(t_i(A_j)\le B_i\), and deletes a chore from \(i\)'s own bundle; deleting from the other bundle would be the goods convention and is inappropriate for chores. Define \(\alpha(A)=\inf\{\alpha\ge1:A\text{ satisfies every such inequality}\}\), with \(+\infty\) if none does. For an explicitly restricted instance class \(\mathcal D\), its guarantee is \(\sup_{I\in\mathcal D}\inf_{A\text{ complete feasible}}\alpha(A)\). Zero comparison disutility can make this infinite; positive processing times do not imply positive disutility. The research task is an axiomatized economically appropriate comparison and its domain-specific guarantee, not an unqualified universal factor.
\subsection*{Variants and assumption changes}
\begin{enumerate}
\item \textbf{Common times (editorial benchmark).} Impose \(t_{ic}=t_c\) and equal budgets. Determine the feasibility--EF1 frontier for additive disutilities, specifying whether chores can be discarded or all must be assigned.
\item \textbf{Heterogeneous times (source agenda, editorial criterion).} Impose a bounded ratio \(t_{ic}/t_{jc}\le R\), positive disutilities and a stated workload-slack condition. Seek finite approximation guarantees, then test whether feasible-subset comparisons are normatively preferable to comparisons of full bundles.
\end{enumerate}
\subsection*{References and source audit}
\begin{enumerate}
\item Dagstuhl §4.1. Supports the stated research area; exact displayed specification is editorial. \url{https://drops.dagstuhl.de/storage/04dagstuhl-reports/volume14/issue09/24401/html/DagRep.14.9.145/DagRep.14.9.145.html}
\end{enumerate}
\begin{enumerate}\raggedright
\item\hypertarget{definitions:source:30007559:1}{} Edith Elkind. 2024. \emph{Fair Division of Chores with Budget Constraints, in Fair Division: Algorithms, Solution Concepts, and Applications}. Seminar report §4.1.
\url{https://drops.dagstuhl.de/storage/04dagstuhl-reports/volume14/issue09/24401/html/DagRep.14.9.145/DagRep.14.9.145.html}.
DOI: \url{https://doi.org/10.4230/DagRep.14.9.145}.
\end{enumerate}

\subsection*{Common conventions: Source mathematical conventions}

These conventions apply to each entry unless explicitly replaced.

\textbf{Models and identification.} A primitive model \(m\) specifies all probability laws, feasible choices, information, equilibrium and measurement rules used by its target. The admissible class \(\mathcal M\) is fixed independently of outcomes. If \(P_O(m)\) is its observable probability law and \(\tau(m)\) the target, its identified set at observable law \(P\) is
\[
 \mathcal I_\tau(P)=\{\tau(m):m\in\mathcal M,\ P_O(m)=P\}.
\]
Point identification means that this set is a singleton. An empty set signals incompatibility of data and assumptions. Coordinate bounds need not describe the joint set. Bounds are sharp only if every retained target is attainable or arbitrarily approachable by compatible models. Finite bounds require explicit restrictions on unobserved quantities; unrestricted confounding, hidden wealth, extrapolation or arbitrary equilibrium selection can yield uninformative sets.

\textbf{Causal interventions.} A potential outcome \(Y_i(p)\) is the outcome under a complete policy regime \(p\), including timing, financing, information and market spillovers. The target population and units are fixed before treatment. With interference, use the complete assignment vector or a justified exposure mapping; individual no-interference notation alone cannot identify an economy-wide intervention. A difference of mean potential outcomes does not require the cross-world joint distribution. Conditioning on post-treatment migration, survival, enrollment or employment changes the estimand and needs separate justification.

\textbf{Statistical statements.} An estimator, confidence set, forecast or test specifies a sampling law, model class and loss/coverage criterion. Uniform \((1-\alpha)\)-coverage means \(\inf_{m\in\mathcal M}\Pr_m\{\tau(m)\in C_n\}\geq1-\alpha\), or an explicitly asymptotic version. The whole parameter space trivially has coverage, so informativeness requires a stated width, risk or power objective. A forecasting improvement is relative to a named benchmark, information set and evaluation population.

\textbf{Equilibrium and optimization.} An equilibrium comprises feasible plans, optimality, beliefs where needed, budget constraints and all market/resource clearing. The primitives or a stated benchmark define each correspondence \(\mathcal E(p;m)\). Nonemptiness is assumed or proved, never inferred just from compact policy sets. A selected-equilibrium target uses a declared selection rule; a robust target quantifies over all permitted equilibria. Use suprema or infima unless attainment is established. An \(\varepsilon\)-optimal policy is feasible and within \(\varepsilon\) of the finite optimum value. Report an empty feasible set separately.

\textbf{Welfare and accounting.} Normative utility, interpersonal normalization, social weights, numeraire, discounting and the use of public receipts are primitives. Behavioral utility need not coincide with the welfare criterion. Household welfare includes ownership income and rents once; adding profits again to compensating variation generally double-counts them. A monetary equivalent is defined by an explicit transfer and price convention, with monotonicity and range conditions. Average effects, mechanism contributions, transfers and welfare losses are different targets. Infinite sums require convergence and dynamic feasible plans require terminal or no-Ponzi conditions.

\textbf{Source versions.} A working-paper date, revision date and journal publication year are distinguished. A DOI for a journal version is not automatically the DOI of an earlier manuscript. Locators refer to the stated version and printed pages where specified. A metadata-only check does not verify a claim about a full-text conclusion. Every entry's source subsection narrows the provenance claim accordingly.
% END ECONOMICS STATEMENT
\end{document}
